% =========================================================
% Proof stubs: Limit Points
% Source: volume-iii/analysis/point-set-topology/notes/limit-points/notes-limit-points.tex
% =========================================================
% Dependency order:
%   1. seq-char-limit     (direct from definition)
%   2. inf-nbhd           (uses seq-char-limit)
%   3. derived-closed     (uses seq-char-limit)
%   4. closed-derived     (uses seq-char-limit + definition of closed)
%   5. closure-formula    (uses derived-closed + closed-derived)
% =========================================================


% ---------------------------------------------------------
\subsection*{Sequential Characterization of Limit Points}
\label{prf:seq-char-limit}
% ---------------------------------------------------------

\begin{remark}[Return]
\hyperref[thm:seq-char-limit]{$\leftarrow$ Back to Theorem (Sequential Characterization of Limit Points) in Notes}
\end{remark}

\begin{theorem}[Sequential Characterization of Limit Points]
Let $(X,d)$ be a metric space, $E \subseteq X$, and $x \in X$.
Then
\[
x \in E'
\iff
\exists (x_n) \subseteq E \text{ with all } x_n \neq x
\text{ and } x_n \to x.
\]
\end{theorem}

\begin{proof}
% TODO
\end{proof}

\begin{remark}[Proof shape]
% TODO: Which direction is harder and why? What construction is the key step?
\end{remark}

\begin{remark}[Dependencies]
% TODO: Definition~\ref{def:limit-point} (limit point), definition of
% metric space convergence, and the Archimedean principle for choosing
% radii $r = 1/n$.
\end{remark}


% ---------------------------------------------------------
\subsection*{Infinite Neighbourhood Characterization}
\label{prf:inf-nbhd}
% ---------------------------------------------------------

\begin{remark}[Return]
\hyperref[cor:inf-nbhd]{$\leftarrow$ Back to Corollary (Infinite Neighbourhood Characterization) in Notes}
\end{remark}

\begin{corollary}[Infinite Neighbourhood Characterization]
Let $(X,d)$ be a metric space and $E \subseteq X$.
Then
\[
x \in E'
\iff
\forall r > 0,\ B_r(x) \cap E
\text{ is infinite}.
\]
\end{corollary}

\begin{proof}
% TODO
\end{proof}

\begin{remark}[Proof shape]
% TODO: Both directions use the Sequential Characterization
% (Theorem~\ref{thm:seq-char-limit}). Which direction needs an
% additional counting argument?
\end{remark}

\begin{remark}[Dependencies]
% TODO: Theorem~\ref{thm:seq-char-limit} (Sequential Characterization).
% The $(\Rightarrow)$ direction also needs that a convergent sequence
% of distinct points produces infinitely many distinct points in every ball.
\end{remark}


% ---------------------------------------------------------
\subsection*{Derived Set is Closed}
\label{prf:derived-closed}
% ---------------------------------------------------------

\begin{remark}[Return]
\hyperref[thm:derived-closed]{$\leftarrow$ Back to Theorem (Derived Set is Closed) in Notes}
\end{remark}

\begin{theorem}[Derived Set is Closed]
Let $(X,d)$ be a metric space and $E \subseteq X$.
Then $E'$ is a closed subset of $X$.
\end{theorem}

\begin{proof}
% TODO
\end{proof}

\begin{remark}[Proof shape]
% TODO: Show $E'$ is closed by showing it is sequentially closed.
% Take any $(x_n) \subseteq E'$ with $x_n \to x$ and show $x \in E'$.
% The triangle inequality and a shrinking-radius argument are the key tools.
\end{remark}

\begin{remark}[Dependencies]
% TODO: Definition~\ref{def:derived-set} (derived set),
% Definition~\ref{def:limit-point} (limit point),
% Theorem~\ref{thm:seq-char-limit} (Sequential Characterization),
% triangle inequality for $d$.
\end{remark}


% ---------------------------------------------------------
\subsection*{Closed Sets and Derived Sets}
\label{prf:closed-derived}
% ---------------------------------------------------------

\begin{remark}[Return]
\hyperref[thm:closed-derived]{$\leftarrow$ Back to Theorem (Closed Sets and Derived Sets) in Notes}
\end{remark}

\begin{theorem}[Closed Sets and Derived Sets]
Let $(X,d)$ be a metric space and $E \subseteq X$.
Then
\[
E \text{ is closed}
\iff
E' \subseteq E.
\]
\end{theorem}

\begin{proof}
% TODO
\end{proof}

\begin{remark}[Proof shape]
% TODO: A biconditional proof. For each direction, unpack the relevant
% sequential definition of closedness and show it is equivalent to
% containment of limit points.
\end{remark}

\begin{remark}[Dependencies]
% TODO: Definition of closed set (sequentially: limits of sequences
% in $E$ remain in $E$), Definition~\ref{def:derived-set},
% Theorem~\ref{thm:seq-char-limit} (Sequential Characterization).
\end{remark}


% ---------------------------------------------------------
\subsection*{Closure Formula}
\label{prf:closure-formula}
% ---------------------------------------------------------

\begin{remark}[Return]
\hyperref[thm:closure-formula]{$\leftarrow$ Back to Theorem (Closure Formula) in Notes}
\end{remark}

\begin{theorem}[Closure Formula]
Let $(X,d)$ be a metric space and $E \subseteq X$.
Then
\[
\overline{E} = E \cup E'.
\]
\end{theorem}

\begin{proof}
% TODO
\end{proof}

\begin{remark}[Proof shape]
% TODO: Show the two sets are equal by double containment.
% Use Theorem~\ref{thm:derived-closed} (derived set is closed) to
% establish that $E \cup E'$ is closed, then appeal to the minimality
% characterisation of closure.
\end{remark}

\begin{remark}[Dependencies]
% TODO: Definition of closure $\overline{E}$ (smallest closed set
% containing $E$), Theorem~\ref{thm:derived-closed} (derived set is
% closed), Theorem~\ref{thm:closed-derived} (closed iff $E'\subseteq E$).
\end{remark}